\scanpage{72}
\mspage{35}
dans ce dernier iso, $P_0$ correspond à la section à l'infini, et $P_1$ à la section $2$ de $\mathbb P^1_{\mathbb Z}$.

D'autre part, on a
\[
(88)\qquad \widehat M_{D_1}\simeq\widehat M_{G_1}\simeq\widehat M_{1,1}
\deleted{\simeq\bigl(\mathbb P^1_{\mathbb Z},\text{ avec signatures }4\text{ et }6\text{ le long des sections }0\text{ et }j_1=2^{\alpha}3^{\beta}\text{ (correspond. à la courbe elliptique d'autom }\mathbb Z/4\mathbb Z\bigr)}
\]
\deleted{dans le dernier iso, \deleted{\unsure{lorsque}} la section à l'infini de $\widehat M_{D_1}$}

Je me rends compte qu'il me faudrait mieux comprendre la structure de $\widehat M_{1,1}$ (et $=$ de $M_{1,1}$), à commencer dans le cadre transcendant (analytique complexe, avec évt. la structure réelle explicite -- i.e. considération de la conjugaison complexe\ldots). J'y reviendrai plus loin.

\uline{Notons déjà}

\uline{Proposition} Si $\dim M_G\geq2$, alors $M_G^\infty$ est \inserted{géom.} connexe \inserted{(fibre par fibre)}.

Dém. Écrivons
\[
(89)\qquad\left\{\begin{aligned}
\widehat M_G&\simeq\prod_{\alpha\in S}\widehat M^{!}_{g_\alpha,\widehat\nu_\alpha}\\
M_G&\simeq\prod_{\alpha\in S}M^{!}_{g_\alpha,\widehat\nu_\alpha}
\end{aligned}\right.
\]
\margin{71}
